% A LaTeX template for EXECUTIVE SUMMARY of the MSc Thesis submissions to 
% Politecnico di Milano (PoliMi) - School of Industrial and Information Engineering
%
% S. Bonetti, A. Gruttadauria, G. Mescolini, A. Zingaro
% e-mail: template-tesi-ingind@polimi.it
%
% Last Revision: October 2021
%
% Copyright 2021 Politecnico di Milano, Italy. NC-BY

\documentclass[11pt,a4paper,twocolumn]{article}

%------------------------------------------------------------------------------
%	REQUIRED PACKAGES AND  CONFIGURATIONS
%------------------------------------------------------------------------------
% PACKAGES FOR TITLES
\usepackage{titlesec}
\usepackage{color}

% PACKAGES FOR LANGUAGE AND FONT
\usepackage[utf8]{inputenc}
\usepackage[english]{babel}
\usepackage[T1]{fontenc} % Font encoding

\usepackage{tikz}% A package for high-quality hand-made figures.
\usetikzlibrary{}
\usetikzlibrary{positioning}
\usetikzlibrary{shapes.symbols}

% PACKAGES FOR IMAGES
\usepackage{graphicx}
\usepackage{transparent} % Enables transparent images
\usepackage{eso-pic} % For the background picture on the title page
\usepackage{subfig} % Numbered and caption subfigures using \subfloat.
\usepackage{tikz} % A package for high-quality hand-made figures.
\usetikzlibrary{positioning, shapes.geometric, arrows.meta, fit}
\graphicspath{{./Images/}} % Directory of the images
\usepackage{caption} % Coloured captions
\usepackage{xcolor} % Coloured captions
\usepackage{amsthm,thmtools,xcolor} % Coloured "Theorem"
\usepackage{float}

% STANDARD MATH PACKAGES
\usepackage{amsmath}
\usepackage{amsthm}
\usepackage{bm}
\usepackage[overload]{empheq}  % For braced-style systems of equations
\usepackage{amssymb}

% PACKAGES FOR TABLES
\usepackage{tabularx}
\usepackage{longtable} % tables that can span several pages
\usepackage{colortbl}

% PACKAGES FOR ALGORITHMS (PSEUDO-CODE)
\usepackage{algorithm}
\usepackage{algorithmic}

% PACKAGES FOR REFERENCES & BIBLIOGRAPHY
\usepackage[colorlinks=true,linkcolor=black,anchorcolor=black,citecolor=black,filecolor=black,menucolor=black,runcolor=black,urlcolor=black]{hyperref} % Adds clickable links at references
\usepackage{cleveref}
\usepackage[square, numbers, sort&compress]{natbib} % Square brackets, citing references with numbers, citations sorted by appearance in the text and compressed
\bibliographystyle{plain} % You may use a different style adapted to your field

% PACKAGES FOR THE APPENDIX
\usepackage{appendix}

% PACKAGES FOR ITEMIZE & ENUMERATES 
\usepackage{enumitem}

% OTHER PACKAGES
\usepackage{amsthm,thmtools,xcolor} % Coloured "Theorem"
\usepackage{comment} % Comment part of code
\usepackage{fancyhdr} % Fancy headers and footers
\usepackage{lipsum} % Insert dummy text
\usepackage{tcolorbox} % Create coloured boxes (e.g. the one for the key-words)
\usepackage{stfloats} % Correct position of the tables
\usepackage{graphicx}   % for resizebox

%-------------------------------------------------------------------------
%	NEW COMMANDS DEFINED
%-------------------------------------------------------------------------
% EXAMPLES OF NEW COMMANDS -> here you see how to define new commands
\newcommand{\bea}{\begin{eqnarray}} % Shortcut for equation arrays
\newcommand{\eea}{\end{eqnarray}}
\newcommand{\e}[1]{\times 10^{#1}}  % Powers of 10 notation
\newcommand{\mathbbm}[1]{\text{\usefont{U}{bbm}{m}{n}#1}} % From mathbbm.sty
\newcommand{\pdev}[2]{\frac{\partial#1}{\partial#2}}
% NB: you can also override some existing commands with the keyword \renewcommand

%----------------------------------------------------------------------------
%	ADD YOUR PACKAGES (be careful of package interaction)
%----------------------------------------------------------------------------


%----------------------------------------------------------------------------
%	ADD YOUR DEFINITIONS AND COMMANDS (be careful of existing commands)
%----------------------------------------------------------------------------


% Do not change Configuration_files/config.tex file unless you really know what you are doing. 
% This file ends the configuration procedures (e.g. customizing commands, definition of new commands)
% Set the geometric layout of the document
\usepackage{geometry}
\geometry{
  top=3cm,
  left = 2.0cm,
  right = 2.0cm,
  bottom=2cm,
  headheight= 2cm,
  headsep= 0cm,
}
\raggedbottom 

% Create color bluePoli (-> manuale grafica coordinata:  https://www.polimi.it/fileadmin/user_upload/il_Politecnico/grafica-coordinata/2015_05_11_46xy_manuale_grafica_coordinata.pdf)
\definecolor{bluePoli}{cmyk}{0.4,0.1,0,0.4}

% Custom theorem environments
\declaretheoremstyle[
  headfont=\color{bluePoli}\normalfont\bfseries,
  bodyfont=\color{black}\normalfont\itshape,
]{colored}

\captionsetup[figure]{labelfont={color=bluePoli}} % Set colour of the captions
\captionsetup[table]{labelfont={color=bluePoli}} % Set colour of the captions
\captionsetup[algorithm]{labelfont={color=bluePoli}} % Set colour of the captions

\theoremstyle{colored}
\newtheorem{theorem}{Theorem}[section]
\newtheorem{proposition}{Proposition}[section]

% Enhances the features of the standard "table" and "tabular" environments.
\newcommand\T{\rule{0pt}{2.6ex}}
\newcommand\B{\rule[-1.2ex]{0pt}{0pt}}

% Algorithm description
\newcounter{algsubstate}
\renewcommand{\thealgsubstate}{\alph{algsubstate}}
\newenvironment{algsubstates}{
    \setcounter{algsubstate}{0}%
    \renewcommand{\STATE}{%
    \stepcounter{algsubstate}%
    \Statex {\small\thealgsubstate:}\space}
    }{}
    
% Custom theorem environment
\newcolumntype{L}[1]{>{\raggedright\let\newline\\\arraybackslash\hspace{0pt}}m{#1}}
\newcolumntype{C}[1]{>{\centering\let\newline\\\arraybackslash\hspace{0pt}}m{#1}}
\newcolumntype{R}[1]{>{\raggedleft\let\newline\\\arraybackslash\hspace{0pt}}m{#1}}

% Custom itemize environment
\setlist[itemize,1]{label=$\bullet$}
\setlist[itemize,2]{label=$\circ$}
\setlist[itemize,3]{label=$-$}
\setlist{nosep}

% Set separation of columns 
\setlength{\columnsep}{30pt}

% Create command for background pic
\newcommand\BackgroundPic{% Adding background picture
	\put(230,358){
		\parbox[b][\paperheight]{\paperwidth}{%
			\vfill
			\centering
			\transparent{0.4}
			\includegraphics[width=0.5\paperwidth]{raggiera_polimi.eps}%
			\vfill
}}}

% Set indentation
\setlength\parindent{0pt}

% Custom title commands
\titleformat{\section}
{\color{bluePoli}\normalfont\Large\bfseries}
{\color{bluePoli}\thesection.}{1em}{}
\titlespacing*{\section}
{0pt}{2ex}{1ex}

\titleformat{\subsection}
{\color{bluePoli}\normalfont\large\bfseries}
{\color{bluePoli}\thesubsection.}{1em}{}
\titlespacing*{\subsection}
{0pt}{2ex}{1ex}

% Custom headers and footers
\pagestyle{fancy}
\fancyhf{}
      
\fancyfoot{}
\fancyfoot[C]{\thepage} % page
\renewcommand{\headrulewidth}{0mm} % headrule width
\renewcommand{\footrulewidth}{0mm} % footrule width

\makeatletter
\patchcmd{\headrule}{\hrule}{\color{black}\hrule}{}{} % headrule
\patchcmd{\footrule}{\hrule}{\color{black}\hrule}{}{} % footrule
\makeatother

% -> Create the header
\chead[C]{
\centering
\begin{tcolorbox}[arc=0pt, boxrule=0pt, colback=bluePoli!60, width=\textwidth, colupper=white]
    \textbf{Executive summary} \hfill \textbf{\author}  
\end{tcolorbox}
}

% Insert here the info that will be displayed into your Title page 
% -> title of your work
\renewcommand{\title}{Explaining Data Drift through Dynamic Clustering Analysis}
% -> author name and surname
\renewcommand{\author}{Mirko Manset}
% -> MSc course
\newcommand{\course}{Computer Science And Engineering}
% -> advisor name and surname
\newcommand{\advisor}{Prof. Marcello Restelli}
% IF AND ONLY IF you need to modify the co-supervisors you also have to modify the file Configuration_files/title_page.tex (ONLY where it is marked)
\newcommand{\firstcoadvisor}{Alessio Russo, Lorenzo Bisi} % insert if any otherwise comment
% \newcommand{\secondcoadvisor}{Alessio Russo} % insert if any otherwise comment
% -> academic year
\newcommand{\YEAR}{2025-26}

%-------------------------------------------------------------------------
%	BEGIN OF YOUR DOCUMENT
%-------------------------------------------------------------------------
\begin{document}

%-----------------------------------------------------------------------------
% TITLE PAGE
%-----------------------------------------------------------------------------
% Do not change Configuration_files/TitlePage.tex (Modify it IF AND ONLY IF you need to add or delete the Co-advisors)
% This file creates the Title Page of the document
% DO NOT REMOVE SPACES BETWEEN LINES!

\twocolumn[{\begin{@twocolumnfalse}

\AddToShipoutPicture*{\BackgroundPic}

\hspace{-0.6cm}\includegraphics[width=0.6\textwidth]{logo_polimi_ing_indinf.eps}

\vspace{-1mm}
\fontsize{0.3cm}{0.5cm}\selectfont \bfseries \textsc{\color{bluePoli} Executive Summary of the Thesis}\\

\vspace{-0.2cm}
\Large{\textbf{\color{bluePoli}{\title}}}\\

\vspace{-0.2cm}
\fontsize{0.3cm}{0.5cm}\selectfont \bfseries \textsc{\color{bluePoli} Laurea Magistrale in \course}\\

\vspace{-0.2cm}
\fontsize{0.3cm}{0.5cm} \selectfont \bfseries Author: \textsc{\textbf{\author}}\\

\vspace{-0.4cm}
\fontsize{0.3cm}{0.5cm}\selectfont \bfseries Advisor: \textsc{\textbf{\advisor}}\\

% if only ONE co-advisor is present:
\vspace{-0.4cm}
\fontsize{0.3cm}{0.5cm}\selectfont \bfseries Co-advisor: \textsc{\textbf{\firstcoadvisor}}\\
% if more than one co-advisors are present:
%\vspace{-0.4cm}
%\fontsize{0.3cm}{0.5cm}\selectfont \bfseries Co-advisors: \textsc{\textbf{\firstcoadvisor}}\textsc{\textbf{\secondcoadvisor}}\\

\vspace{-0.4cm}
\fontsize{0.3cm}{0.5cm}\selectfont \bfseries Academic year: \textsc{\textbf{\YEAR}}

\small \normalfont

\vspace{11pt}

\centerline{\rule{1.0\textwidth}{0.4pt}}

\vspace{15pt}
\end{@twocolumnfalse}}]

\thispagestyle{plain} % In order to not show the header in the first page

\section{Introduction}\label{ch:introduction}

Machine Learning (ML) is increasingly applied in environments where data arrive
continuously in a streaming manner. Such streaming data often exhibit evolving
statistical properties, a phenomenon known as \emph{data drift}, which poses
challenges for traditional static analysis methods.

Dynamic clustering addresses this issue by incrementally grouping incoming data
and tracking how clusters evolve over time. Unlike conventional approaches, it
focuses on explaining the structural transformations induced by drift rather
than merely detecting it. This provides a human-interpretable view of pattern
evolution, enabling analysts to understand the formation, splitting, merging,
and recurrence of clusters.

By continuously analyzing temporal dynamics, dynamic clustering turns raw,
non-stationary data into actionable insights. It supports informed
decision-making in real-time applications by revealing how and why clusters
change, offering a transparent understanding of evolving data streams.


\section{Preliminaries}\label{sec:preliminaries}

This section introduces foundational concepts essential for understanding
clustering in streaming environments and the challenges of evolving data
distributions.

\subsection{Clustering}\label{subsec:clustering}

Clustering partitions a dataset $D = \{\mathbf{x}_1, \dots, \mathbf{x}_N\}$
into $K$ clusters by optimizing an objective $\mathcal{J}(C; D)$ that
encourages high similarity within clusters and low similarity across clusters.
Classical methods include K-means, which assigns points to nearest centroids,
and Gaussian Mixture Models (GMMs), which model data as mixtures of Gaussian
distributions with parameters $\{\pi_j, \boldsymbol{\mu}_j,
    \mathbf{\Sigma}_j\}$ and provide soft assignments through posterior
probabilities.

\subsection{Streaming Learning}\label{subsec:streaming_learning}

In streaming environments, data arrives continuously as an unbounded sequence.
Streaming algorithms incrementally update a model $\mathcal{M}_t$ with each new
point $\mathbf{x}_t$:
\[
    \mathcal{M}_{t+1} = \mathcal{F}(\mathcal{M}_t, \mathbf{x}_t),
\]
balancing stability and adaptability under bounded memory and computational
constraints.

\subsection{Data Drift}\label{subsec:data_drift}

\textbf{Data drift}~\cite{drift_adaptation_survey} refers to changes in statistical properties over time, occurring when $p_t(X, y) \neq p_\tau(X, y)$ between times $t$ and $\tau$. \emph{Real drift} involves changes in the conditional distribution $p(y \mid X)$, while \emph{virtual drift} affects the marginal distribution $p(X)$ with stable $p(y \mid X)$.

Detection strategies monitor model performance, compare
distributions, or combine multiple signals. Adaptation mechanisms include
\emph{blind adaptation} (continuous updates) and \emph{informed adaptation}
(updates triggered by detected drift), employing retraining, incremental
updates, or ensemble strategies to maintain predictive accuracy under evolving
conditions.

\subsection{Drift Explainability}\label{subsec:drift_explainability}

Drift explainability identifies the causes and features responsible for
distributional changes, clarifying why model performance degrades over time.
\emph{Feature importance analysis} compares contributions $\phi_i$ between
reference and production models, with large discrepancies $\Delta \phi_i =
    |\phi_i^{\text{ref}} - \phi_i^{\text{prod}}|$ revealing influential
drift-driving variables. The \emph{Classifier Two-Sample Test} (C2ST) treats
reference and production datasets as separate classes, training a classifier to
distinguish them, feature importances from this classifier highlight dimensions
contributing most to distributional change.

Explainable detection methods ensure model updates are necessary and
effectively targeted, enhancing trustworthiness and reliability of adaptive
systems in dynamic environments.

\section{Problem Formulation}\label{ch:problem_formulation}

This thesis addresses the problem of clustering in evolving data streams, where
the underlying distributions are non-stationary and subject to drift over time.
Traditional clustering methods assume access to fixed datasets and static
distributions, making them unsuitable for environments where data arrives
continuously and patterns evolve. The aim is to establish a framework that can
process data incrementally while also tracking and interpreting how clusters
change across time.

\subsection{Streaming Clustering}\label{sec:prob_streaming_clustering}

A data stream can be described as a potentially infinite sequence of
$d$-dimensional observations
$\mathcal{X}=\{\mathbf{x}^t\}_{t=1}^{\infty}$. The task of streaming clustering
is to incrementally partition these observations into coherent groups while
meeting strict constraints on memory and computation. At each time $t$, the
algorithm produces a clustering $C^t=\{C_1^t,\dots,C_{k_t}^t\}$, where the
number of clusters $k_t$ and their composition may vary as new data arrives.

Unlike static clustering, streaming clustering must adapt to non-stationary
conditions. Distributions often evolve gradually rather than abruptly, and
patterns may recur over time, reflecting realistic phenomena such as seasonal
effects or user behavior. A useful abstraction is that of \emph{piecewise
regular drift}, where the stream consists of consecutive phases within which
distributional change is smooth, while transitions between phases are bounded.
This perspective avoids both unrealistic assumptions of perfect stability and
the intractability of arbitrary change. Streaming clustering algorithms must
therefore balance two goals: producing stable partitions within phases while
remaining responsive to evolving structure.

\subsection{Tracking Clusters}\label{sec:prob_tracking_clusters}

As data distributions shift, the structure of clusters also changes. The
clustering $C^t$ at time $t$ may differ substantially from $C^\tau$ at a later
time $\tau$, requiring mechanisms to model and track transitions. Typical transitions
include the appearance of new clusters, the disappearance of old ones, splits
and merges, and the survival of clusters that persist with limited change. These
events provide a structural description of how patterns in the data evolve.

To formalize transitions, clusters at different times are compared by measuring
their degree of \emph{overlap}. Depending on the representation, overlap can be defined
by shared elements, similarity of centroids, or proximity of cluster boundaries.
If the overlap exceeds a given threshold, two clusters are considered related.
This principle allows consistent mapping of clusters across time and supports
the identification of cluster transitions.

Beyond \emph{external transitions}, which describe how clusters relate across time,
\emph{internal transitions} capture changes within a single surviving cluster. These
include expansion or shrinkage in size, variation in density, and shifts in the
centroid. While such internal dynamics offer useful insights, the primary focus
in this work is on external transitions, which provide a clearer view of the
structural evolution of the clustering model.

Tracking clusters in dynamic streams presents several challenges. Distributional
drift may be gradual or abrupt, making it difficult to decide when two clusters
should be linked. Noise and outliers may obscure genuine transitions, while
overlapping clusters can create ambiguity. Furthermore, the absence of ground
truth in streaming settings complicates evaluation, requiring surrogate measures
or application-driven validation.

\subsection{Dynamic Clustering}\label{sec:dynamic_clustering}

Dynamic clustering integrates streaming clustering with temporal cluster
tracking, offering a unified framework for understanding the evolution of data
distributions. The objective is not only to efficiently cluster data as it
arrives, but also to interpret how clusters emerge, evolve, and dissolve over
time. This dual perspective provides richer insights than static methods,
especially in domains where understanding dynamics is as valuable as achieving
accurate predictions.

Several analytical goals motivate this framework. First, by observing the
history of clusters, it becomes possible to explain when and how drift occurs
and to link these changes to external factors. Second, monitoring the lifecycle
of clusters supports concept evolution analysis, where patterns emerge,
stabilize, and eventually disappear. Third, sudden transitions, such as abrupt
merges or splits, can serve as indicators of anomalies or significant events in
the underlying system. Finally, dynamic clustering supports interpretability by
enabling visualization of cluster trajectories, timelines, or trees, helping
experts make sense of complex temporal patterns.

In summary, dynamic clustering treats clustering as an evolving process rather
than a one-off computation. It unites the efficiency of online algorithms with
the interpretability required to analyze changes in non-stationary
environments. This formulation provides the foundation for adaptive and
trustworthy clustering systems that can operate effectively under continuous
change.



\section{Method}\label{method}

This section presents the proposed framework for dynamic clustering in evolving
data streams. The method is designed to capture both short-term structural
changes and long-term recurrence patterns by extending cluster transition
modeling, introducing new similarity measures, and integrating them into a
modular clustering architecture.

\subsection{New Transition Types}
This work extends traditional cluster evolution events by introducing three new
recurrence-aware transitions: \textbf{reappearance} (a previously disappeared
cluster re-emerges), \textbf{remerge} (previously split clusters recombine),
and \textbf{resplit} (previously merged clusters split again).

Cluster evolution is modeled using a graph-based representation, where nodes
correspond to clusters at different timestamps and edges indicate overlap
relationships. This structure enables interpretable identification of both
classical transitions (appearance, disappearance, split, merge) and the newly
introduced recurrence transitions, allowing a more realistic modeling of
evolving data streams over extended periods.

\subsection{Overlapping Scores}
To quantify cluster similarity and guide transition detection, two overlapping
scores are introduced.

\paragraph{Spherical Overlapping Score (SOS).}
For approximately spherical clusters $C_i$ and $C_j$, such as those from
K-means, with centers $\boldsymbol{\mu}_i, \boldsymbol{\mu}_j$ and radii $r_i,
    r_j$, the SOS is defined as:
\[
    \text{SOS}(C_i, C_j) = 2^{- \frac{||\boldsymbol{\mu}_i - \boldsymbol{\mu}_j||_2}{r_i + r_j}}.
\]
This value ranges between 0 and 1, decays smoothly with distance, and naturally
defines a threshold at 0.5 to indicate potential evolutionary continuity.

\paragraph{Effective Overlapping Score (EffectiveOS).}
For clusters modeled as Gaussian distributions with means $\boldsymbol{\mu}_i,
    \boldsymbol{\mu}_j$ and covariances $\Sigma_i, \Sigma_j$, EffectiveOS
generalizes SOS by employing the Mahalanobis distance:

{\scriptsize
\begin{align}
    \text{EffectiveOS}(C_i, C_j)                & = 2^{- \frac{d_M(\boldsymbol{\mu}_i, \boldsymbol{\mu}_j)}{\text{EffSize}_i + \text{EffSize}_j}},                    \\
    d_M(\boldsymbol{\mu}_i, \boldsymbol{\mu}_j) & = \sqrt{(\boldsymbol{\mu}_i - \boldsymbol{\mu}_j)^\top \Sigma_{ij}^{-1} (\boldsymbol{\mu}_i - \boldsymbol{\mu}_j)}, \\
    \Sigma_{ij}                                 & = \frac{\Sigma_i + \Sigma_j}{2}.
\end{align}}

Here, $\text{EffSize}_i$ is a confidence-aware effective size based on
principal variances:
\[
    \text{EffSize}_i = \sqrt{\frac{1}{d} \sum_{\ell=1}^{d} (\chi^2_\alpha \sqrt{\lambda_{\ell,i}})^2},
\]
where $\lambda_{\ell,i}$ are the eigenvalues of $\Sigma_i$ and $\chi^2_\alpha$
is the chi-squared critical value at confidence $\alpha$. EffectiveOS remains
bounded in $(0,1]$ while capturing directional variability and correlation in
ellipsoidal clusters.

\subsection{Dynamic Clustering Framework}
The proposed framework is modular, consisting of five integrated components.

\paragraph{Online phase (micro-clustering).}
Streaming data is summarized into compact statistical representations in real
time. In this work, \textit{CluStream}~\cite{clustream} is employed to ensure
efficiency under memory and latency constraints while retaining essential
structure.

\paragraph{Triggering strategy.}  
A strategy determines when to perform macro-clustering, based either on fixed
time intervals or on data-driven signals that indicate significant
distributional change. For simplicity, and to remain coherent with the original
\textit{CluStream} design, we adopt fixed intervals as the triggering mechanism
in this work.

\paragraph{Offline phase (macro-clustering).}
Micro-clusters are aggregated into higher-level groups. Gaussian Mixture Models
provide flexible representations in which the number of clusters can adapt
dynamically to capture complex evolving structures.

\paragraph{History module.}
Minimal descriptors of past clusters are maintained to track long-term
dynamics. This enables the detection of recurrence patterns such as
\emph{reappearance}, \emph{remerge}, and \emph{resplit}.

\paragraph{Tracking module.}
Macro-clusters across different time steps are compared using similarity
measures such as SOS or EffectiveOS. This facilitates the identification of
structural changes and the reconstruction of historical trajectories.

The framework produces a comprehensive report, including graph-based
visualizations of cluster transitions, spatial trajectories of macro-cluster
evolution, and representative data samples.

\begin{figure}[H]
    \centering
    \includegraphics[width=1\linewidth]{Images/my_images/architecture.png}
    \caption{Proposed architecture for the dynamic clustering framework.}
\end{figure}


\section{Experiments}\label{sec:experiments}

This section summarizes the evaluation of the proposed dynamic clustering
algorithm on two datasets: a typographical drift dataset derived from the 20
Newsgroups corpus and a real image dataset of fruits. These experiments highlight the
algorithm's ability to adapt to evolving data distributions for both text and
image data.

\subsection{Typographical Drift Dataset}\label{subsec:typographical_drift}

To assess robustness under linguistic noise, a subset of the 20
Newsgroups~\cite{20newsgroups} dataset was used, consisting of four categories:
\emph{rec.autos}, \emph{comp.graphics}, \emph{sci.med}, and
\emph{rec.sport.baseball}. Initially, clean documents were sampled to form a
reference distribution. Gradually, typographical noise was introduced through
character-level modifications such as insertions, deletions, replacements, and
swaps, simulating progressive degradation of text quality.

Each document was embedded into a 384-dimensional semantic vector using the
\texttt{all-MiniLM-L6-v2} transformer model~\cite{sentence-transformers}, and
reduced to 32 dimensions using UMAP~\cite{umap}. The algorithm was configured
with $q=200$ micro-clusters, a streaming window $\Delta=100$, a
macro-clustering interval $I=100$, an overlapping threshold $\varepsilon=0.5$
and a confidence level $\alpha=0.9$.

\begin{table}[ht]
    \centering
    \scriptsize % small font
    \renewcommand{\arraystretch}{0.9} % reduces row height
    \setlength{\tabcolsep}{4pt} % reduces space between columns
    \begin{tabularx}{\linewidth}{c|l|X}
        \hline
        \textbf{Cluster} & \textbf{Interpretation} & \textbf{Text Sample} \\
        \hline
        0 & sci.med & \emph{Ihf antne has any informatiin on thi deficiency would veyr ggreatly apprecaite a response hereq or preferably by Emaul. All I know at tihs opin is a defjciencg can caue smyoglovin to be releasfd,a dn in times of stresz ad hgh ambivent temperatursy could caussez grenal failure. x} \\
        \hline
        1 & rec.sport.baseball & \emph{Bobby Bonilla supposedly use the word 'faggot' when he got mad at that author in the clubhouse. Should he be banned from baseball for a year like Schott?} \\
        \hline
        3 & comp.graphics & \emph{Ihf antne has a ljst if cimpanies oing dagta isualizatoing (software or ahdwzre)I woild liek toheapzrf rom tyem. Thank.s - - rs --} \\
        \hline
        2 & rec.autos & \emph{Cup holders (driving is an importantant enough undertaking) Cellular phones and mobile fax machines (see above) Vanity mirrors on the driver's side. Ashtrays (smokers seem to think it's just fine to use the road) Fake convertible roofs and vinyl roofs. Any gold trim.} \\
        \hline
        6 & Typos & \emph{Thhe qystion is nkt wgethfe your adio wigll e stolen. Tghe question isw en your ryaeio will bestolen.} \\
        \hline
    \end{tabularx}
    \caption{Final representative samples of text typo clusters.}
\end{table}


\begin{figure}[H]
    \centering
    \includegraphics[width=1\linewidth]{my_images/experiment_results/text_typos/graph.png}
    \caption{Graphical evolution of Typographical Drift clusters over time.}
\end{figure}

Over the course of the experiment, the algorithm tracked the evolution of
clusters in response to increasing noise. Initially, the drift manifested as
subtle splits from existing clusters, with merging and re-splitting observed as
the text gradually degraded. When the text became heavily corrupted, fully
distinct clusters emerged, reflecting pronounced shifts in the data
distribution. These results demonstrate the algorithm's sensitivity to both
gradual and substantial changes, capturing complex temporal dynamics in textual
data.

% \begin{figure}[H]
%     \centering
%     \includegraphics[width=1\linewidth]{my_images/experiment_results/text_typos/selected_spatial.png}
%     \caption{Spatial evolution of Typographical Drift clusters showing temporal progression.}
% \end{figure}

\subsection{Fruits Dataset}\label{subsec:fruits}

The Fruits dataset consists of 2,420 RGB images of tangerines, resized to $100
    \times 100$ pixels. Background segmentation was applied using
GrabCut~\cite{grabcut}, and high-dimensional embeddings were extracted using a
pretrained ResNet-50~\cite{resnet} network. The embeddings were reduced to 64
dimensions with UMAP~\cite{umap}. The first 200 images served as reference
data, and the remaining 2,220 images were progressively ingested.
Hyperparameters were set to $q=200$, $\Delta=50$, $I=50$, $\varepsilon=0.5$ and
$\alpha=0.9$.

\begin{figure}[H]
    \centering
    \includegraphics[width=1\linewidth]{my_images/experiment_results/fruits/selected_rsamples.png}
    \caption{Representative Samples of clusters of Fruit Images.}
\end{figure}

\begin{figure}[H]
    \centering
    \includegraphics[width=1\linewidth]{my_images/experiment_results/fruits/selected_spatial.png}
    \caption{Spatial evolution of clusters of Fruit Images showing temporal progression.}
\end{figure}

\begin{figure}[H]
    \centering
    \includegraphics[width=1\linewidth]{my_images/experiment_results/fruits/graph.png}
    \caption{Graphical evolution of Fruit Images clusters over time.}
\end{figure}

The algorithm successfully captured the emergence of new tangerine types and
gradual visual changes over time. Clusters corresponded to distinct
characteristics such as color (green, yellow, orange) and shape (round or
oval). While minor splits were occasionally missed due to sensitivity
thresholds, major transitions, including the appearance of orange tangerines,
were clearly detected.


\section{Conclusions}

This thesis explored dynamic clustering in non-stationary data streams, where
traditional static methods cannot adapt to evolving distributions and recurring
structures. A modular and adaptive framework was developed to maintain
temporally coherent and interpretable clusters, supported by a typology of
transitions, \emph{survival}, \emph{absorption}, \emph{splitting}, and long-term
patterns such as \emph{reappearance}, \emph{remerge}, and \emph{resplit}. Novel
overlap metrics, including the \emph{Spherical Overlapping Score (SOS)} and the
\emph{Effective Overlapping Score (EffectiveOS)}, provide geometry-aware criteria to robustly
detect cluster transitions, enhancing interpretability and robustness under
gradual drift or recurring patterns.

The framework supports incremental updates, memory-efficient summarization,
visualization, and cluster-level reporting. Experiments demonstrated its
effectiveness in capturing evolving cluster structures over time and providing
human-interpretable insights into data dynamics.

Future developments include extending the framework to density-based clustering
algorithms (e.g., DBSCAN, HDBSCAN) by designing new overlap metrics suited to
non-centroid clusters. The modular architecture allows integration of
alternative clustering methods and distance functions, as well as advanced
triggering strategies informed by drift detection. Another promising direction
is forecasting cluster transitions, leveraging the temporal metadata produced
by the framework to anticipate structural changes and enable timely
adaptation in dynamic data streams.


\bibliography{bibliography.bib} 

\end{document}